\section{Esperimento}
\label{sec:experiments}

\subsection{Protocollo e separazione del costo offline}

Sono confrontati A (DFS con branch-and-bound), B (ricerca con ordinamento
predittivo) e C (ricerca predittiva con regolatore e soglia $P=100$).
Ogni esecuzione parte dal deposito, con percorso $p=(0)$, sei consegne
da visitare e costo accumulato nullo. I tre metodi ricevono la stessa
matrice e lo stesso budget $B\in\{50,100,200,500\}$ di stati estratti.
Ogni prova riparte da zero: i budget maggiori non riprendono una pila
lasciata dalle prove precedenti.

Il modello principale è quello addestrato sui seed $42$--$49$.
I seed $50$ e $51$, esclusi dall'addestramento dei pesi, sono stati usati
anche per esaminare e modificare la politica di regolazione: vanno quindi
considerati casi di sviluppo della politica. Sono state successivamente
generate due nuove istanze, con seed $52$ e $53$, mantenendo fissi modello,
soglia e regola. Queste istanze sono usate soltanto per la valutazione,
senza generare ulteriori esempi di allenamento o riaddestrare il modello.

La generazione dei target e l'addestramento sono operazioni offline.
Anche l'ottimo di riferimento è calcolato separatamente, una volta per
istanza, con B\&B e budget $100\,000$. La ricerca di riferimento termina
in tutti e quattro i casi. I costi ottimi sono rispettivamente
$288{,}252135$, $309{,}557853$, $259{,}896129$ e $349{,}580289$ per i
seed $50$, $51$, $52$ e $53$. Servono soltanto per valutare i risultati:
non inizializzano la migliore soluzione e non guidano i metodi online.

Il budget misura estrazioni, comprendendo stati terminali e stati poi
potati, non soltanto espansioni. B e C richiedono anche valutazioni del
modello; C esegue inoltre scansioni della pila durante la diversificazione.
Lo stesso numero di estrazioni non implica uguale tempo di esecuzione.
Non sono riportate misure temporali, quindi il confronto riguarda la
qualità a parità di estrazioni e non dimostra un vantaggio di velocità.

\subsection{Misure di qualità e risultati A/B/C}

Per ogni esecuzione si registrano il miglior costo totale $G_B$, il
percorso, il completamento della ricerca, le statistiche di frontiera e
potatura e, per C, il numero di diversificazioni. La qualità è misurata
anche tramite il divario percentuale dall'ottimo:
\[
    \operatorname{gap}_{\%}=100\,\frac{G_B-G^*}{G^*}.
\]
Il divario viene calcolato soltanto se è disponibile un giro e il
riferimento è certificato. Il codice gestisce il caso di ottimo nullo
evitando la divisione per zero.

\begin{table}[htbp]
\centering
\small
\caption{Costi totali su sei consegne, incluso il ritorno al deposito.
A: baseline; B: predittiva; C: predittiva regolata.
$D_C$ conta le estrazioni di diversificazione.}
\label{tab:search-results}
\begin{tabular}{rrrrrr}
\toprule
Seed & Budget & Costo A & Costo B & Costo C & $D_C$\\
\midrule
50 & 50  & $384{,}226$ & $309{,}735$ & $309{,}735$ & 0\\
50 & 100 & $380{,}162$ & $309{,}735$ & $309{,}735$ & 0\\
50 & 200 & $309{,}735$ & $288{,}252$ & $288{,}252$ & 0\\
50 & 500 & $288{,}252$ & $288{,}252$ & $288{,}252$ & 3\\
\midrule
51 & 50  & $336{,}449$ & $369{,}651$ & $369{,}651$ & 0\\
51 & 100 & $331{,}273$ & $364{,}901$ & $364{,}901$ & 0\\
51 & 200 & $331{,}273$ & $331{,}273$ & $309{,}558$ & 1\\
51 & 500 & $309{,}558$ & $309{,}558$ & $309{,}558$ & 4\\
\midrule
52 & 50  & $296{,}682$ & $259{,}896$ & $259{,}896$ & 0\\
52 & 100 & $296{,}682$ & $259{,}896$ & $259{,}896$ & 0\\
52 & 200 & $278{,}501$ & $259{,}896$ & $259{,}896$ & 1\\
52 & 500 & $259{,}896$ & $259{,}896$ & $259{,}896$ & 4\\
\midrule
53 & 50  & $425{,}883$ & $403{,}590$ & $403{,}590$ & 0\\
53 & 100 & $425{,}883$ & $381{,}082$ & $381{,}082$ & 0\\
53 & 200 & $394{,}277$ & $354{,}992$ & $354{,}992$ & 0\\
53 & 500 & $349{,}580$ & $354{,}992$ & $349{,}580$ & 2\\
\bottomrule
\end{tabular}
\end{table}

Le sedici righe rappresentano $48$ esecuzioni. In tutte, il numero di
stati estratti è pari al budget e la pila non è vuota all'arresto:
\texttt{is\_complete=False}. Sono comunque disponibili giri completi.
Un costo uguale all'ottimo di riferimento identifica un giro ottimo,
entro la precisione numerica, ma non una certificazione ottenuta dal
metodo online: trovare l'ottimo e completare la ricerca sono eventi distinti.

Sul seed $51$ a budget $200$, A e B hanno gap $7{,}01\%$, mentre C
raggiunge l'ottimo con una diversificazione. Sul seed $52$, B trova
l'ottimo già entro $50$ estrazioni e C non migliora ulteriormente il costo.
Sul seed $53$, a budget $200$ i gap di A e B sono rispettivamente
$12{,}79\%$ e $1{,}55\%$; C non è ancora intervenuto. A budget $500$,
B resta all'$1{,}55\%$, mentre A e C raggiungono l'ottimo; C effettua
due diversificazioni. In questo caso il guadagno della regolazione è
rispetto a B, mentre il costo finale coincide con quello di A.

\subsection{Controllo separato su istanze didattiche}

Il modello didattico della Sezione~\ref{sec:teaching-data} è stato
verificato separatamente su mappe con quattro consegne, usando budget
$\{5,10,20,50\}$. Sulla mappa \texttt{test\_irregolare}, l'ottimo di
riferimento è $15{,}512156$. A budget $5$, A trova costo $19{,}196$
(gap $23{,}75\%$) e B costo $17{,}997$ (gap $16{,}02\%$). Entro $10$
estrazioni B raggiunge l'ottimo; A resta a $19{,}196$ anche a budget
$20$ e raggiunge l'ottimo entro $50$. C coincide con B: con soglia $100$
non interviene in queste prove. Questo controllo illustra l'effetto
dell'ordinamento appreso, ma non costituisce una valutazione della
regolazione né sostituisce l'esperimento principale su sei consegne.
